% Source coverage: physical PDF pp. 35--37 (printed pp. 12--14), ending before Section 15.4.
% Equations: (15.3.1)--(15.3.10), plus intervening unnumbered displays.
% Footnotes: none. Citations: Ref. 8. Uncertainties: none.

\subsection{Field Equations and Conservation Laws}
\label{sec:15.3}

Using Eq.~\eqref{eq:15.2.9} for the matrix $g_{\alpha\beta}$ in
Eq.~\eqref{eq:15.2.3}, the full Lagrangian density is
\begin{equation}
 \mathcal L=-\frac14F_{\alpha\mu\nu}F_\alpha{}^{\mu\nu}
 +\mathcal L_M(\psi,D_\mu\psi) ,
 \tag{15.3.1}\label{eq:15.3.1}
\end{equation}
where in the absence of gauge fields $\mathcal L_M(\psi,\partial_\mu\psi)$ would
be the `matter' Lagrangian density. We could, in principle, include a dependence
of $\mathcal L_M$ on $F_{\alpha\mu\nu}$ as well as higher covariant derivatives
$D_\nu D_\mu\psi$, $D_\lambda F_{\alpha\mu\nu}$, etc., but we exclude these
non-renormalizable terms here for the same reason as in electrodynamics: as
discussed in Section 12.3, such terms would be highly suppressed at ordinary
energies by negative powers of some very large mass. For this reason the standard
model of the weak, electromagnetic and strong interactions has a Lagrangian of
the general form \eqref{eq:15.3.1}.

The equations of motion of the gauge field are here
\begin{align*}
 \partial_\mu\frac{\partial\mathcal L}{\partial(\partial_\mu A_{\alpha\nu})}
 &=-\partial_\mu F_\alpha{}^{\mu\nu}
 =\frac{\partial\mathcal L}{\partial A_{\alpha\nu}}\\
 &=-F_\gamma{}^{\nu\mu}C_{\gamma\alpha\beta}A_{\beta\mu}
 -i\frac{\partial\mathcal L_M}{\partial D_\nu\psi}t_\alpha\psi
\end{align*}
and so
\begin{equation}
 \partial_\mu F_\alpha{}^{\mu\nu}=-\mathcal J_\alpha{}^\nu ,
 \tag{15.3.2}\label{eq:15.3.2}
\end{equation}
where $\mathcal J_\alpha{}^\nu$ is the current:
\begin{equation}
 \mathcal J_\alpha{}^\nu\equiv
 -F_\gamma{}^{\nu\mu}C_{\gamma\alpha\beta}A_{\beta\mu}
 -i\frac{\partial\mathcal L_M}{\partial D_\nu\psi}t_\alpha\psi .
 \tag{15.3.3}\label{eq:15.3.3}
\end{equation}
The current $\mathcal J_\alpha{}^\nu$ is conserved in the ordinary sense
\begin{equation}
 \partial_\nu\mathcal J_\alpha{}^\nu=0 ,
 \tag{15.3.4}\label{eq:15.3.4}
\end{equation}
as can be seen either from the Euler--Lagrange equations for $\psi$ and the
invariance equivalent \eqref{eq:15.2.2} or, more easily, directly from the field
equations \eqref{eq:15.3.2}.

The derivatives in Eqs.~\eqref{eq:15.3.2} and \eqref{eq:15.3.4} are ordinary
derivatives, not the gauge-covariant derivatives $D_\nu$, so the gauge invariance
of these equations is somewhat obscure. It can be made manifest by rewriting
Eq.~\eqref{eq:15.3.2} in terms of the gauge-covariant derivative of the field
strength
\begin{align}
 D_\lambda F_\alpha{}^{\mu\nu}
 &\equiv\partial_\lambda F_\alpha{}^{\mu\nu}
 -i(t^A_\beta)_{\alpha\gamma}A_{\beta\lambda}F_\gamma{}^{\mu\nu}\notag\\
 &=\partial_\lambda F_\alpha{}^{\mu\nu}
 -C_{\alpha\gamma\beta}A_{\beta\lambda}F_\gamma{}^{\mu\nu} .
 \tag{15.3.5}\label{eq:15.3.5}
\end{align}
Then Eq.~\eqref{eq:15.3.2} reads
\begin{equation}
 D_\mu F_\alpha{}^{\mu\nu}=-J_\alpha{}^\nu ,
 \tag{15.3.6}\label{eq:15.3.6}
\end{equation}
where $J_\alpha{}^\nu$ is the current of the matter fields alone
\begin{equation}
 J_\alpha{}^\nu\equiv-i\frac{\partial\mathcal L_M}{\partial D_\nu\psi}
 t_\alpha\psi .
 \tag{15.3.7}\label{eq:15.3.7}
\end{equation}
This is gauge-covariant, if $\mathcal L_M$ is gauge-invariant. Also, by operating
on Eq.~\eqref{eq:15.3.6} with $D_\nu$, using the commutation relation
\[
 [D_\nu,D_\mu]F_\alpha{}^{\mu\nu}
 =-i(t^A_\gamma)_{\alpha\beta}F_{\gamma\nu\mu}F_\beta{}^{\mu\nu}
 =-C_{\gamma\alpha\beta}F_{\gamma\nu\mu}F_\beta{}^{\mu\nu} ,
\]
we see that $J_\alpha{}^\nu$ satisfies a gauge-covariant conservation law
\begin{equation}
 D_\nu J_\alpha{}^\nu=0 ,
 \tag{15.3.8}\label{eq:15.3.8}
\end{equation}
rather than the ordinary conservation law \eqref{eq:15.3.4} obeyed by the full
current $\mathcal J_\alpha{}^\nu$. Also, it is straightforward (using
Eq.~\eqref{eq:15.1.5}) to derive the identities:
\begin{equation}
 D_\mu F_{\alpha\nu\lambda}+D_\nu F_{\alpha\lambda\mu}
 +D_\lambda F_{\alpha\mu\nu}=0 ,
 \tag{15.3.9}\label{eq:15.3.9}
\end{equation}
which hold whether or not the gauge fields satisfy the field equations.

These results serve to underscore the profound analogy mentioned in Section 15.1
between non-Abelian gauge theories and general relativity. In general relativity
there is a matter energy-momentum tensor $T^\nu{}_\mu$, analogous to
$J_\alpha{}^\mu$, which satisfies a generally covariant conservation law
$T^\nu{}_{\mu;\nu}=0$, and stands on the right-hand side of the Einstein field
equations in their generally covariant form,
$R^\nu{}_\mu-\frac12\delta^\nu{}_\mu R=-8\pi G T^\nu{}_\mu$. However,
$T^\nu{}_\mu$ is not conserved in the ordinary sense:
$\partial_\nu T^\nu{}_\mu$ does not vanish. On the other hand, moving the
non-linear terms on the left-hand side of the Einstein equation to the right-hand
side gives a field equation~\hyperref[ch15-ref-8]{[8]}
\[
 \left(R^\nu{}_\mu-\frac12\delta^\nu{}_\mu R\right)_{\mathrm{LINEAR}}
 =-8\pi G\tau^\nu{}_\mu ,
\]
where $\tau^\nu{}_\mu$ is the non-tensor
\[
 \tau^\nu{}_\mu\equiv T^\nu{}_\mu
 +\frac{1}{8\pi G}
 \left(R^\nu{}_\mu-\frac12\delta^\nu{}_\mu R\right)_{\mathrm{NONLINEAR}} ,
\]
analogous to $\mathcal J_\alpha{}^\nu$. Like $\mathcal J_\alpha{}^\nu$,
$\tau^\nu{}_\mu$ is conserved in the ordinary sense
\[
 \partial_\nu\tau^\nu{}_\mu=0
\]
and may be regarded as the current of energy and momentum:
\[
 P_\mu=\int\tau^0{}_\mu\,d^3x .
\]
It contains a purely gravitational term, because gravitational fields carry
energy and momentum; without this term, $\tau^\nu{}_\mu$ could not be conserved.
Similarly, $\mathcal J_\alpha{}^\nu$ contains a gauge-field term (the first term
on the right in Eq.~\eqref{eq:15.3.3}) because for non-Abelian groups (those with
$C^\gamma{}_{\alpha\beta}\ne0$) the gauge fields carry the quantum numbers with
which they interact. Because $\mathcal J_\alpha{}^\nu$ is conserved in the
ordinary sense, it can be regarded as the current of these quantum numbers, with
the symmetry generators given by the time-independent quantities
\begin{equation}
 T_\alpha=\int\mathcal J_\alpha{}^0\,d^3x .
 \tag{15.3.10}\label{eq:15.3.10}
\end{equation}
(Also, the homogeneous equations \eqref{eq:15.3.9} involve covariant derivatives,
just as do the Bianchi identities of general relativity.) In contrast, none of
these complications arises in quantum electrodynamics, because photons do not
carry the quantum number, electric charge, with which they interact.
